\documentclass[11pt]{article}
\usepackage[margin=25mm]{geometry}
\usepackage{fontspec}
\setmainfont{Latin Modern Roman}
\usepackage{amsmath,amssymb,mathtools}
\usepackage{unicode-math}
\setmathfont{Latin Modern Math}
\usepackage{xurl,hyperref}
\hypersetup{colorlinks=true,urlcolor=blue}
\setcounter{secnumdepth}{0}
\setlength{\parindent}{0pt}
\setlength{\parskip}{5pt}
\setlength{\emergencystretch}{3em}
\newcommand{\sref}[2]{\hyperlink{source:#1:#2}{[S#2]}}
\newcommand{\cataloguescope}[1]{\paragraph{Recorded target.}#1}
\newcommand{\R}{\mathbb R}
\newcommand{\C}{\mathbb C}
\newcommand{\N}{\mathbb N}
\newcommand{\Z}{\mathbb Z}
\newcommand{\Q}{\mathbb Q}
\newcommand{\F}{\mathbb F}
\newcommand{\eps}{\varepsilon}
\begin{document}
\section{Large-Scale Mechanisms Driving Atmospheric Monsoons}\label{statement}

\subsection{Definitions and mathematical statement}

\paragraph{Shared notation.}$\mathbb N=\{1,2,\ldots\}$; the standard number fields have their usual meanings. All additional conventions are stated locally.


\paragraph{Definitions and conventions.} A monsoon is a seasonal reversal of large-scale atmospheric circulation and associated precipitation patterns. Near-surface wind reversal and cross-equatorial overturning are characteristic features in the source discussion. Thermal contrasts force circulation, while transport of angular momentum by atmospheric eddies can alter the large-scale balance. The report compares mechanisms using simplified atmosphere models.

\paragraph{Source scope and qualifications.} This is a mechanism question for Earth's atmosphere, not a single quantified theorem with a fixed model and observable. The report develops a proposed explanation; it does not imply that a universal necessary-and-sufficient monsoon condition is already defined.

\paragraph{Catalogue mathematical target.}
\paragraph{Catalogue wording retained for traceability.}
What drives large-scale monsoons in the Earth's atmosphere?

\paragraph{Saved literature context.} The supplied catalogue's assessment dated 2026-08-21 records: The report frames monsoon forcing as an incompletely understood research question; later atmospheric-science work refines competing mechanisms but does not provide a universally accepted complete explanation. The assessment is retained as dated provenance; it is not a new comprehensive literature or open-status certification.
\subsection{Short English statement}
Explain the dynamical mechanisms driving the seasonal large-scale monsoon circulation.
\subsection{Research attempt}
\paragraph{Model, date and investigated scope.}
OpenAI GPT-6 Astra (configured model \texttt{gpt-6-astra}), ultra reasoning; generated 11 October 2026. Solves a seasonal two-variable forcing model and derives its amplitude, lag and strict sign-reversal condition.

\paragraph{Sources and literature check.}
The complete supplied definition, including source qualifications and the dated literature assessment, was reviewed. The calculation below is an elementary restricted attempt motivated by [S2]; it does not claim novelty or a fresh literature-status determination.

\paragraph{Context and investigated scope.} The complete supplied definition, including source qualifications and the dated literature assessment, was reviewed. The calculation below is an elementary restricted attempt motivated by [S2]; it does not claim novelty or a fresh literature-status determination.

\paragraph{A seasonal thermal-forcing calculation.} Let $T$ denote a hemispheric temperature contrast and $M$ a signed overturning amplitude in the reduced linear model
\[
 C\dot T+kT=Q\cos(\omega t),\qquad \dot M+rM=aT,
\]
with $C,k,r,\omega>0$. The unique long-time periodic solution is the real part of
\[
 T(t)=\frac{Qe^{i\omega t}}{k+i\omega C},\qquad
 M(t)=\frac{aQe^{i\omega t}}{(k+i\omega C)(r+i\omega)}.
\]
Substitution verifies both equations; all transients decay as combinations of $e^{-kt/C}$ and $e^{-rt}$ (or a polynomial times that exponential if the rates coincide).

\paragraph{Reversal criterion in this restricted model.} The overturning amplitude is $A=|aQ|/[\sqrt{k^2+\omega^2C^2}\sqrt{r^2+\omega^2}]$. With a background mean $M_0$, the flow $M_0+A\cos(\omega t-\varphi)$ changes sign exactly when $A>|M_0|$; equality only touches zero. Thermal inertia produces phase lag $\arctan(\omega C/k)$, with additional dynamical lag $\arctan(\omega/r)$. Thus identical heating amplitudes can give different reversal behavior through heat capacity and damping.

\paragraph{What is not explained.} The coefficients and mean offset were imposed, and the model contains no moisture, precipitation, angular-momentum constraints or eddy transport. It therefore gives a testable reduced transfer relation, not a derivation of Earth's monsoons. The first gap is identifying and validating the momentum/thermal closure in the source atmosphere regime, including whether nonlinear transitions dominate the smooth response just calculated. No universal monsoon criterion is claimed.

\paragraph{Verification and adversarial audit.}
Checked complex response, coincident transient rates, background offset and the equality case in reversal; coefficients remain explicit model assumptions. Model self-audit only; no independent human proof checking. Prompts A and B were applied to this canonical record and its qualified source scope. Checked complex response, coincident transient rates, background offset and the equality case in reversal; coefficients remain explicit model assumptions. Retained unresolved status and zero original-target progress; A justified atmosphere closure including moisture, eddies and angular-momentum transport is missing.

Coordinator checked the two linear periodic responses, decaying transients and strict sign-reversal amplitude criterion. This additional check was model-assisted, not independent human verification.

\paragraph{Tools, process and contributions.}
Python 3 was used for file preparation and structural checks; no numerical experiment or theorem prover was used. Contribution prompt A was followed by prompt B. GPT-6 Astra with ultra reasoning drafted the research and performed its stated self-audit; any additional coordinator contribution is identified in the verification and audit. A GPT-6-based Codex coordinator assembled the manuscript and checked repository integration; its exact variant and reasoning configuration were not exposed. No human mathematical contribution, independent human proof check or formal proof certification is supplied. Complete derivations and the consulted source references are included here; private reasoning transcripts are not published.

\paragraph{Final status and first remaining gap.}
\textbf{UNRESOLVED}. A justified atmosphere closure including moisture, eddies and angular-momentum transport is missing.

\section{Completion Estimate}
\noindent\textbf{COMPLETION: 0\%}
The calculation is an auxiliary or restricted attempt and establishes no newly verified portion of the full source target. A justified atmosphere closure including moisture, eddies and angular-momentum transport is missing.

\subsection{Sources}
\begin{itemize}
\item[\hypertarget{source:statement:1}{S1}] ULAM.AI and the UnsolvedMath Contributors, UnsolvedMath, catalogue entry. Catalogue curation under CC BY 4.0; definitions and English exposition added locally.
\url{https://huggingface.co/datasets/ulamai/UnsolvedMath}.
\item[\hypertarget{source:statement:2}{S2}] Mathematical Theory and Modelling in Atmosphere-Ocean-Science (2010), Oberwolfach Reports, DOI 10.4171/owr/2010/34, pp. 2027--2099; official report source and mathematical context.
\url{https://doi.org/10.4171/owr/2010/34}.
\item[\hypertarget{source:statement:3}{S3}] Official EMS Press PDF of the same Oberwolfach report; PDF pages 11--12.
\url{https://ems.press/content/serial-article-files/46292}.
\end{itemize}
\label{end:statement}
\end{document}
